\documentclass[../../../include/open-logic-section]{subfiles}

\begin{document}

\olfileid{sth}{replacement}{limofsize}
\olsection{د اندازې محدوديت}

ښايي د تعويض د «دننۍ» توجيې تر ټولو عامه هڅه د لاندې تصور له لارې وي:
\begin{enumerate}
	\item[] \limofsize. هر څيزونه يو سټ جوړوي، په دې شرط چې ډېر زيات نۀ وي.
\end{enumerate}
دا اصل به تعويض سمدستي توجيه کړي۔ په پای کښې، د تعويض په مټ جوړ شوے
هر سټ له هغه سټ څخه لوے کېدے نۀ شي چې ترې جوړ شوے وي۔ په کره توګه:
فرض کړئ د تعويض په کارولو سره
$\funimage{\tau}{A} = \Setabs{\tau(x)}{x \in A}$ جوړوئ؛ نو
$\cardle{\funimage{\tau}{A}}{A}$۔ په دې توګه، که د $A$ !!{element}s
د سټ جوړولو دپاره تر اندازې ډېر نۀ وو، نو د هغو انځورونه هم د
$\funimage{\tau}{A}$ د جوړولو دپاره تر اندازې ډېر نۀ دي۔

د تعويض د توجيې دپاره د \limofsize{} په کارولو کښې ښکاره ستونزه دا ده
چې تر اوسه مو د \limofsize په څېر هېڅ اصل \emph{نۀ} دے اېښے۔ سربېره
پر دې، کله چې مو د سټ د تجمعي-تکراري تصور کيسه په
\crefrange{sth:story::chap}{sth:z::chap} کښې وکړه، په هغې کښې
د \limofsize لوري ته هېڅ \emph{اشاره} هم نۀ وه۔ همدا وجه ده چې بولوس
يو وخت وليکل: «ښايي دې پايلې ته ورسېږو چې د سټ تيورۍ تر شا لږ تر لږه
دوه فکرونه پراته دي» \citeyearpar[p.~19]{Boolos1989}۔ له يوې خوا، د سټ
د تجمعي-تکراري تصور شاوخوا مفکورې د~$\Z$ توجيه کول غواړي۔ له بلې خوا،
\limofsize{} د تعويض توجيه کول غواړي۔

خو ستونزه يوازې دا نۀ ده چې تر اوسه د \limofsize په اړه \emph{چپ}
وو۔ اصلي ستونزه دا ده چې \limofsize{}، په همدغه بڼه چې بيان شو، د سټ
له تجمعي-تکراري تصور سره چندان سم نۀ ښکاري۔ ځکه په دې اصل کښې په
\emph{پړاوونو} کښې د سټونو د جوړېدو تصور هېڅ نۀ يادېږي۔

د دغو «دوو فکرونو» تاريخ ته په کتو دا حيرانوونکې نۀ ده؛ يعنې د سټ
تجمعي-تکراري تصور او \limofsize دوه بېلې لارې دي چې د سټ تيورۍ
تناقضونه پرې بندېږي۔ تجمعي-تکراري تصور د رسل تناقض په لنډه توګه داسې
بندوي: \emph{موږ بايد د رسل د سټ د شتون تمه نۀ واے کړې، ځکه په هېڅ
پړاو کښې به «جوړ» شوے نۀ واے}۔ په مقابل کښې، \limofsize{} د رسل سټ
په دې وينا ردوي: \emph{موږ بايد د رسل د سټ د شتون تمه نۀ واے کړې،
ځکه هغه به تر اندازې ډېر لوے واے}۔

نو په دې بڼه ئې په ښکاره وايو: د تناقضونو د ځواب په توګه،
\limofsize{} \emph{ډېرې} زياتې توجيې ته اړ دے۔ د بېلګې په توګه، د رسل
د تناقض دا بڼه وګورئ: \emph{هېڅ داسې پګ سپے نشته چې عيناً هغه او يوازې
هغه پګ سپي بوي کړي چې خپل ځان نۀ بوي کوي} (\olref[sth][story][rus]{sec}
وګورئ)۔ که پوښتنه وکړئ «ولې داسې پګ سپے نشته؟»، دا ښه ځواب نۀ دے چې
هغه به ډېر زيات پګ سپي بوي کول غواړي۔ نو د رسل د سټ د نۀ شتون دا به
ولې ښه بداهتي توضيح وي چې د شتون دپاره به «تر اندازې ډېر لوے» واے؟

په لنډه توګه، که د \limofsize د «بداهتي» انګېزې په اړه لږ حيران ياست،
دا د پوهېدو وړ ده۔

\end{document}
